\documentclass[11pt]{article}

\usepackage[T1]{fontenc}
\usepackage{amsmath}
\usepackage{amssymb}
\usepackage{amsthm}
\usepackage{booktabs}
\usepackage[hidelinks]{hyperref}
\usepackage[margin=1in]{geometry}

\newtheorem{theorem}{Theorem}
\newtheorem{lemma}[theorem]{Lemma}
\newtheorem{corollary}[theorem]{Corollary}
\newcommand{\Rcal}{\mathcal{R}}

\title{Prime-Order Automorphism Exclusions for Ramsey $(5,5;43)$ Graphs}
\author{Ruturaj R Raval\\Independent Researcher\\
ORCID: 0000-0003-4930-8981}
\date{September 7, 2026}

\begin{document}
\maketitle

\begin{abstract}
Let $\Rcal(5,5;43)$ denote the graphs on 43 vertices with clique number and
independence number at most four. Elementary fixed-point, degree, and
incidence arguments exclude 16 prime-order automorphism cycle types. A dated
public-source audit found 12 of those types previously uncovered or
unfinished. We also exclude the intervening order-3 type, $3^6 1^{25}$,
using four orbit-CNF branches whose compressed binary DRAT proofs are
retained and independently checked. A second reduction excludes
$3^8 1^{19}$ using elementary incidence arguments, an exact structural
slack identity, and ten checked orbit-CNF branches. The combined original
result closes 14 cycle types that the audited public records listed as
uncovered or unfinished. The result does not determine $R(5,5)$ or change
the interval $43 \le R(5,5) \le 46$.
\end{abstract}

\section{Context}

The diagonal Ramsey number $R(5,5)$ is the least $n$ such that every graph on
$n$ vertices contains a clique or an independent set of size five. Exoo
proved $R(5,5)\ge 43$ in 1989 \cite{Exoo1989}. McKay and Radziszowski proved
$R(5,5)\le49$ in 1997 \cite{McKayRadziszowski1997}. Angeltveit and McKay
improved the upper bound to 48 in 2018 \cite{AngeltveitMcKay2018} and to 46
in 2025 \cite{AngeltveitMcKay2025}. Thus the current interval is
\[
43 \le R(5,5) \le 46.
\]

This paper studies a finite structural subproblem. If a graph in
$\Rcal(5,5;43)$ exists, which prime-order automorphisms can it admit? For a
prime $p$, an automorphism with $c$ moved $p$-cycles and $43-pc$ fixed
vertices has cycle type
\[
p^c 1^{43-pc}.
\]

\section{Elementary exclusions}

\begin{lemma}[Degree interval]
\label{lem:degree}
Every vertex of a graph $G\in\Rcal(5,5;43)$ has degree between 18 and 24.
\end{lemma}

\begin{proof}
For a vertex $v$, the graph induced by $N(v)$ contains neither a clique of
size four nor an independent set of size five. Since $R(4,5)=25$
\cite{McKayRadziszowski1995}, we have $d(v)\le24$. The nonneighbors of $v$
contain neither a clique of size five nor an independent set of size four.
Again using $R(5,4)=25$, we obtain $42-d(v)\le24$, hence $d(v)\ge18$.
\end{proof}

\begin{lemma}[Fixed-point bound]
\label{lem:fixed}
Let $G\in\Rcal(5,5;43)$ admit an automorphism of prime order $p\ge5$ with a
moved orbit. Then the automorphism has at most 26 fixed vertices.
\end{lemma}

\begin{proof}
Let $C$ be a moved $p$-cycle. The induced graph $G[C]$ is neither complete
nor empty, since either alternative contains a homogeneous set of size five.
Thus $G[C]$ contains an edge and a nonedge.

Every fixed vertex is adjacent to all or none of $C$. Let $S$ be the fixed
vertices complete to $C$, and let $T$ be the fixed vertices anticomplete to
$C$. If $S$ contained a triangle, that triangle together with an edge of
$C$ would form a clique of size five. The graph $G[S]$ is therefore
triangle-free and has no independent set of size five. Since $R(3,5)=14$
\cite{GreenwoodGleason1955}, $|S|\le13$.

If $T$ contained an independent triple, that triple together with a nonedge
of $C$ would form an independent set of size five. The graph $G[T]$ has no
clique of size five, so $R(5,3)=14$ gives $|T|\le13$. Hence the number of
fixed vertices is at most $|S|+|T|\le26$.
\end{proof}

\begin{theorem}[Elementary cycle-type exclusions]
\label{thm:elementary}
No graph in $\Rcal(5,5;43)$ admits a prime-order automorphism of any of the
following 16 cycle types:
\[
\begin{aligned}
&2^c1^{43-2c} &&(1\le c\le3),\\
&3^c1^{43-3c} &&(1\le c\le5),\\
&3^7 1^{22},\\
&5^c1^{43-5c} &&(1\le c\le3),\\
&7^c1^{43-7c} &&(1\le c\le2),\\
&11^1 1^{32},\qquad 13^1 1^{30}.&
\end{aligned}
\]
\end{theorem}

\begin{proof}
For $p\ge5$, Lemma~\ref{lem:fixed} excludes every listed type because it has
more than 26 fixed vertices.

Now let $p=3$. Each moved 3-cycle induces either a triangle or an independent
triple. Suppose first that a moved cycle is a triangle. Its fixed common
neighbors form an independent set and therefore have size at most four. A
vertex on the cycle has degree at most
\[
2+3(c-1)+4=3c+3.
\]
This contradicts Lemma~\ref{lem:degree} when $c\le4$.

If a moved cycle is an independent triple, its fixed common nonneighbors form
a clique and therefore have size at most four. A vertex on the cycle has at
least
\[
(43-3c)-4=39-3c
\]
fixed neighbors. This exceeds 24 when $c\le4$.

At $c=5$, both bounds meet the degree interval exactly. Equality for a
triangle cycle forces it to be complete to all other moved cycles. Equality
for an independent cycle forces it to be anticomplete to all other moved
cycles. A triangle cycle and an independent cycle cannot coexist. If all
five cycles are triangles, the 15 moved vertices form a clique. If all five
are independent, they form an independent set. Both cases are impossible.

It remains to treat seven moved 3-cycles. For distinct moved cycles $C_i$
and $C_j$, the nine possible cross edges split into three invariant perfect
matchings. Let $w_{ij}\in\{0,1,2,3\}$ be the number present and set
\[
q_i=\sum_{j\ne i}w_{ij}.
\]
If $C_i$ is a triangle, at most four fixed vertices are complete to it, so
the degree lower bound gives $q_i\ge12$. Two triangle cycles satisfy
$w_{ij}\le2$, since complete linkage creates a clique of size at least five.
If $C_i$ is independent, at least 18 fixed vertices are complete to it, so
the degree upper bound gives $q_i\le6$. Two independent cycles satisfy
$w_{ij}\ge1$, since empty linkage creates an independent set of size six.

Suppose $t$ moved cycles are triangles and $s=7-t$ are independent. Summing
$q_i$ over the triangle cycles gives
\[
12t\le \sum_{i:\,C_i\text{ triangle}}q_i.
\]
Triangle-triangle links contribute at most $4\binom{t}{2}$. Each independent
cycle uses at least $s-1$ of its total weight on the other independent
cycles, and therefore contributes at most $6-(s-1)=t$ to triangle cycles.
Hence
\[
12t\le4\binom{t}{2}+st=t(t+5),
\]
which is impossible for $1\le t\le6$. All seven cycles have the same type.
After complementation, assume they are all independent.

Each of the six intercycle weights at a moved cycle is at least one, while
$q_i\le6$. Thus every intercycle weight is one and $q_i=6$. The degree upper
bound and the 18-fixed-neighbor lower bound then force each moved cycle to
have exactly 18 fixed neighbors and four fixed nonneighbors.

For a fixed vertex $x$, let $Z_x$ be the set of moved cycles anticomplete to
$x$. Counting fixed-cycle nonneighbor incidences by moved cycle gives
\[
\sum_{x\text{ fixed}}|Z_x|=7\cdot4=28. \tag{1}
\]
If fixed vertices $x$ and $y$ are adjacent, their common neighborhood
contains neither a triangle nor an independent set of size five. Therefore
$R(3,5)=14$ gives $|N(x)\cap N(y)|\le13$. Every moved cycle outside
$Z_x\cup Z_y$ contributes all three vertices to that common neighborhood,
so
\[
3(7-|Z_x\cup Z_y|)\le13,
\]
and hence $|Z_x\cup Z_y|\ge3$.

The fixed vertices with $|Z_x|\le1$ are consequently pairwise nonadjacent
and number at most four. Every other fixed vertex has $|Z_x|\ge2$, which
would force
\[
\sum_{x\text{ fixed}}|Z_x|\ge2(22-4)=36.
\]
This contradicts (1), so the type $3^7 1^{22}$ is impossible.

Finally let $p=2$. A moved pair is an edge or a nonedge. For an edge pair,
the fixed common-neighbor graph is triangle-free and independent-5-free, so
it has at most 13 vertices. A vertex in the pair has degree at most
\[
1+2(c-1)+13=2c+12.
\]
For a nonedge pair, the fixed common-nonneighbor graph is clique-5-free and
independent-triple-free, so it has at most 13 vertices. A vertex in the pair
has degree at least
\[
(43-2c)-13=30-2c.
\]
These bounds exclude $c=1,2$. At $c=3$, equality forces every edge pair to
be complete to the other moved vertices and every nonedge pair to be
anticomplete to them. Mixed status is inconsistent. Uniform edge status
gives a clique on six vertices, while uniform nonedge status gives an
independent set on six vertices.
\end{proof}

\section{\texorpdfstring{The cycle type $3^6 1^{25}$}
{The cycle type 3\string^6 1\string^25}}

\subsection{Orbit-CNF encoding}

Fix the permutation
\[
(0\,1\,2)(3\,4\,5)\cdots(15\,16\,17)
\]
with vertices 18 through 42 fixed. One Boolean variable represents each
orbit of unordered vertex pairs. There are 501 edge-orbit variables.

For every 5-subset $X$, the encoding adds one clause forbidding all edges of
$X$ and one clause forbidding all nonedges of $X$. Duplicate orbit variables
within a clause are removed. Exact sequential counters impose
$18\le d(v)\le24$ on one representative of each moved vertex orbit and on
every fixed vertex. Repeated literals in a degree sum retain edge-orbit
multiplicity.

\subsection{Complete branch cover}

Choose fixed vertex 18 as a root. Its adjacency to each moved 3-cycle is
constant. Permuting the six moved cycles maps every root pattern to one in
which the adjacent cycles form a prefix of length $k$. Complementation maps
$k$ to $6-k$ and preserves membership in $\Rcal(5,5;43)$. It is therefore
enough to solve $k=0,1,2,3$. An executable audit enumerates all $2^6=64$
patterns and confirms this reduction.

Each branch has 67,709 variables and 910,918 clauses. An independent artifact
auditor checks the DIMACS header, clause count, literal range, file digest,
common formula body, independently reconstructed root variables, and the six
branch-defining unit clauses.

\subsection{Checked certificates}

Kissat reported UNSAT for all four branches. The proof files are compressed
binary DRAT traces. Each was decompressed and independently checked with
\texttt{drat-trim}. Table~\ref{tab:certificates} gives digest prefixes; the
repository manifest records every full SHA-256 value.

\begin{table}[h]
\centering
\begin{tabular}{rrrrrr}
\toprule
$k$ & Variables & Clauses & CNF SHA prefix & Proof SHA prefix & Check s\\
\midrule
0 & 67709 & 910918 & \texttt{ac56dd4e6633} & \texttt{c4365b739438} & 6.043\\
1 & 67709 & 910918 & \texttt{2a92ff17114f} & \texttt{e23f985a7dca} & 6.590\\
2 & 67709 & 910918 & \texttt{5543127a0f17} & \texttt{6c570981d5c2} & 6.298\\
3 & 67709 & 910918 & \texttt{e8a8e8453f45} & \texttt{ecc372452863} & 7.596\\
\bottomrule
\end{tabular}
\caption{Retained certificates for the four root branches.}
\label{tab:certificates}
\end{table}

\begin{theorem}[Certified order-3 exclusion]
\label{thm:c6}
No graph in $\Rcal(5,5;43)$ admits an automorphism of cycle type
$3^6 1^{25}$.
\end{theorem}

\begin{proof}
Any such graph satisfies the orbit-CNF for exactly one root pattern. Cycle
relabeling and complementation map that pattern to one of the four encoded
branches. Every branch is UNSAT by a checked DRAT proof.
\end{proof}

\section{\texorpdfstring{The cycle type $3^8 1^{19}$}
{The cycle type 3\string^8 1\string^19}}

\subsection{Internal types and the elementary extremes}

Let $C_1,\ldots,C_8$ be the moved 3-cycles. Each $C_i$ induces either a
triangle or an independent triple. Let $t$ be the number of triangle cycles
and put $s=8-t$.

For an independent cycle $I_j$, let $B_j$ be the fixed vertices
anticomplete to $I_j$. The set $B_j$ is a clique, since two nonadjacent
vertices of $B_j$ together with $I_j$ would form an independent set of size
five. Hence $|B_j|\le4$. Dually, the fixed vertices complete to a triangle
cycle form an independent set and also number at most four.

\begin{lemma}[Extreme internal-type counts]
\label{lem:c8-extremes}
The values $t=0,1,7,8$ are impossible.
\end{lemma}

\begin{proof}
First suppose $t=1$, so there are seven independent cycles. For a fixed
vertex $x$, let $z_x$ be the number of independent cycles anticomplete to
$x$. Counting incidences with the sets $B_j$ gives
\[
\sum_x z_x\le7\cdot4=28. \tag{2}
\]
If fixed vertices $x$ and $y$ both satisfy $z_x,z_y\le1$, then they are
complete to at least five common independent cycles. Were $xy$ an edge,
their common neighborhood would contain at least 15 vertices. The common
neighborhood of an adjacent pair contains neither a triangle nor an
independent set of size five, so $R(3,5)=14$ gives a contradiction. Thus the
fixed vertices with $z_x\le1$ form an independent set and number at most
four. Every other fixed vertex has $z_x\ge2$, forcing
\[
\sum_x z_x\ge2(19-4)=30,
\]
contrary to (2).

Now suppose $t=0$. The same argument shows that the fixed vertices with
$z_x\le1$ form an independent set of size at most four, while
\[
\sum_x z_x\le8\cdot4=32. \tag{3}
\]
Let $n_0$ and $n_1$ count fixed vertices with $z_x=0$ and $z_x=1$.
Since $n_0+n_1\le4$, if $n_0\le1$ then
\[
\sum_x z_x\ge n_1+2(19-n_0-n_1)\ge33,
\]
contrary to (3). Hence $n_0\ge2$.

A fixed vertex with $z_x=0$ is adjacent to all 24 moved vertices. The degree
upper bound therefore makes it anticomplete to every fixed vertex. Choose
two such vertices. The other 17 fixed vertices contain neither a clique of
size five nor an independent triple, since an independent triple together
with the chosen two vertices would form an independent set of size five.
This contradicts $R(5,3)=14$.

Complementation sends $t$ to $8-t$, so it also excludes $t=7,8$.
\end{proof}

\subsection{The exact slack identity}

It remains to consider $t=2,3,4$, since complementation covers $t=5,6$.
Write the triangle cycles as $T_1,\ldots,T_t$ and the independent cycles as
$I_1,\ldots,I_s$.

For $T_i$, let $a_i$ be its number of fixed neighbors, put
$\alpha_i=4-a_i$, and write its degree as $18+u_i$. For $I_j$, let $b_j$ be
its number of fixed nonneighbors, put $\beta_j=4-b_j$, and write its degree
as $24-v_j$. All four quantities $\alpha_i,\beta_j,u_i,v_j$ are
nonnegative.

The invariant edges between two moved cycles split into three perfect
matchings. A triangle-triangle link has weight at most two; write
$D_{ik}=2-w_{ik}\ge0$. An independent-independent link has weight at least
one; write $B_{j\ell}=w_{j\ell}-1\ge0$. Let $m_{ij}$ be the mixed weight
between $T_i$ and $I_j$, and define the mixed row and column sums
\[
R_i=\sum_jm_{ij},\qquad C_j=\sum_im_{ij}.
\]

The degree equation at $T_i$ is
\[
18+u_i=2+(4-\alpha_i)
 +\sum_{k\ne i}(2-D_{ik})+R_i,
\]
and therefore
\[
R_i=14-2t+\alpha_i+u_i+\sum_{k\ne i}D_{ik}. \tag{4}
\]
The degree equation at $I_j$ is
\[
24-v_j=(15+\beta_j)
 +\sum_{\ell\ne j}(1+B_{j\ell})+C_j,
\]
and therefore
\[
C_j=t+2-\beta_j-v_j-\sum_{\ell\ne j}B_{j\ell}. \tag{5}
\]
Summing (4) and (5), then using $\sum_iR_i=\sum_jC_j$, gives the exact
identity
\[
\sum_i\alpha_i+\sum_j\beta_j+\sum_iu_i+\sum_jv_j
 +2\sum_{i<k}D_{ik}+2\sum_{j<\ell}B_{j\ell}
 =(t-4)^2. \tag{6}
\]

Every term in (6) is a nonnegative integer. Consequently:
\[
\begin{aligned}
\sum_i a_i+\sum_j b_j&\ge32-(t-4)^2,\\
\sum_{i<k}D_{ik}+\sum_{j<\ell}B_{j\ell}
 &\le\left\lfloor\frac{(t-4)^2}{2}\right\rfloor,\\
2t(7-t)\le\sum_{i,j}m_{ij}&\le(8-t)(t+2),\\
R_i&\ge14-2t,\qquad C_j\le t+2.
\end{aligned} \tag{7}
\]
The encoder also uses the row upper and column lower bounds obtained by
subtracting the other rows or columns from the two mixed-total bounds.

\subsection{Complete low-exception branch cover}

Each moved cycle has at most four fixed-vertex exceptions, so the total
number of exception incidences is at most 32. Among 19 fixed vertices there
is therefore a root with at most one exception. Relabel triangle cycles and
independent cycles separately. For $t=2,3$, the normalized root branches are
\[
(t,0,0),\qquad(t,1,0),\qquad(t,0,1),
\]
where the last two coordinates count adjacent triangle exceptions and
nonadjacent independent exceptions. For $t=4$, complementation identifies
$(4,0,1)$ with $(4,1,0)$.

An executable audit enumerates all $2^8$ internal-type patterns and the nine
root patterns with zero or one exception. Of the 2,304 configurations, 162
belong to Lemma~\ref{lem:c8-extremes}. The remaining 2,142 configurations
reduce exactly as shown in Table~\ref{tab:c8-cover}.

\begin{table}[h]
\centering
\begin{tabular}{rr}
\toprule
Branch & Configurations\\
\midrule
$(2,0,0)$ & 56\\
$(2,1,0)$ & 112\\
$(2,0,1)$ & 336\\
$(3,0,0)$ & 112\\
$(3,1,0)$ & 336\\
$(3,0,1)$ & 560\\
$(4,0,0)$ & 70\\
$(4,1,0)$ & 560\\
\bottomrule
\end{tabular}
\caption{Normalized low-exception root branches for $3^8 1^{19}$.}
\label{tab:c8-cover}
\end{table}

\subsection{The zero-slack case}

When $t=4$, the right side of (6) is zero. Thus every term on the left is
zero. In particular, each triangle cycle has exactly four fixed neighbors
and degree 18, each independent cycle has exactly four fixed nonneighbors
and degree 24, every triangle-triangle weight is two, every
independent-independent weight is one, and every mixed row and column has
sum six.

\begin{lemma}[Mixed-matrix classification]
\label{lem:c8-matrix}
Every mixed weight is one or two. Up to independent permutations of the four
triangle cycles and four independent cycles, the mixed-weight matrix is one
of
\[
\begin{pmatrix}
1&1&2&2\\
1&1&2&2\\
2&2&1&1\\
2&2&1&1
\end{pmatrix},
\qquad
\begin{pmatrix}
1&1&2&2\\
1&2&1&2\\
2&1&2&1\\
2&2&1&1
\end{pmatrix}. \tag{8}
\]
\end{lemma}

\begin{proof}
Let $A_i$ be the four fixed neighbors of triangle cycle $T_i$, and let
$B_j$ be the four fixed nonneighbors of independent cycle $I_j$. If
$m_{ij}=3$, then every vertex of $A_i$ must be nonadjacent to $I_j$, since
otherwise $T_i$, that fixed vertex, and one vertex of $I_j$ form a
5-clique. Hence $A_i\subseteq B_j$, and equality follows from their sizes.
But $A_i$ is independent while $B_j$ is a clique, a contradiction.
Complementation excludes $m_{ij}=0$.

Every entry is therefore one or two. Since each row and column sums to six,
subtracting one from each entry gives the adjacency matrix of a 2-regular
simple bipartite graph on four plus four vertices. Its cycle decomposition
is either one 8-cycle or two 4-cycles, giving the two matrices in (8). A
direct enumeration finds 90 labeled matrices and exactly these two classes,
even when the first triangle row is distinguished by the root.
\end{proof}

The strengthened formulas also use three immediate consequences. Two fixed
vertices sharing a triangle exception are nonadjacent, since otherwise they
join that triangle to form a 5-clique. Two fixed vertices sharing an
independent exception are adjacent, since otherwise they join that
independent triple to form an independent set of size five. Hence
$|A_i\cap B_j|\le1$. Independent rotations of the seven nonroot moved
cycles fix seven matching phases on a spanning star. In the $(4,1,0)$
branch, a fixed vertex with zero exceptions would instead be selected as the
$(4,0,0)$ root, so zero-exception fixed signatures are excluded to make the
two root branches disjoint.

\subsection{Checked certificates}

The six $t=2,3$ branches use the canonical structural formula. The two
$t=4$ root branches split by Lemma~\ref{lem:c8-matrix}, giving four
strengthened formulas. Table~\ref{tab:c8-formulas} records their dimensions.

\begin{table}[h]
\centering
\begin{tabular}{lrrr}
\toprule
Family & Branches & Variables & Clauses\\
\midrule
$t=2$ canonical & 3 & 76,728 & 880,702\\
$t=3$ canonical & 3 & 76,440 & 879,868\\
$t=4$, root $(4,0,0)$ & 2 & 78,411 & 886,323\\
$t=4$, root $(4,1,0)$ & 2 & 78,411 & 886,342\\
\bottomrule
\end{tabular}
\caption{Retained certificate formulas for $3^8 1^{19}$.}
\label{tab:c8-formulas}
\end{table}

An independent reference implementation reconstructs every clause in all
ten formulas. A separate audit verifies the branch arithmetic, the
zero-slack matrix classification, DIMACS integrity, metadata, and SHA-256
digests. Kissat reported every branch UNSAT. Each source proof was reduced
to a compact core and both the extraction and the retained compact proof
were independently checked with \texttt{drat-trim}.

\begin{theorem}[Certified order-3 exclusion]
\label{thm:c8}
No graph in $\Rcal(5,5;43)$ admits an automorphism of cycle type
$3^8 1^{19}$.
\end{theorem}

\begin{proof}
Lemma~\ref{lem:c8-extremes} excludes four internal-type counts.
Complementation reduces the others to $t=2,3,4$. The exhaustive
low-exception audit reduces every remaining configuration to one of the six
canonical or four matrix-split branches. Each branch is UNSAT by a checked
DRAT proof.
\end{proof}

\section{Novelty boundary and limitations}

A targeted audit of public literature and repositories through September 7,
2026 found the following 12 types uncovered or unfinished before
Theorem~\ref{thm:elementary}:
\[
2^c1^{43-2c}\ (1\le c\le3),\qquad
3^c1^{43-3c}\ (1\le c\le5),\qquad
3^7 1^{22},\qquad
5^c1^{43-5c}\ (1\le c\le3).
\]
The same audit found $3^6 1^{25}$ listed as uncovered or unfinished in the
closest public prescribed-automorphism records \cite{Wustep2026,Kriebel2026}.
It also found $3^8 1^{19}$ pending in the current case list and uncovered in
the independent coverage ledger.
The order-7, order-11, and order-13 consequences of
Theorem~\ref{thm:elementary} are elementary reproofs of cases already covered
computationally.

The $3^6 1^{25}$, $3^7 1^{22}$, and $3^8 1^{19}$ cases had been prepared as
pending finite computations in the closest public campaign, so the
contribution is their resolution rather than their first attempted
formulation.

The combined original result therefore excludes 14 previously uncovered or
unfinished cycle types. This is a structural restriction, not a new Ramsey
bound. It does not address graphs with trivial automorphism group and does
not imply that a hypothetical graph is asymmetric. Many other prime-order
and composite-order automorphism types remain untreated.

The computational exclusions depend on the correctness of the orbit encoder,
the proof checker, and the executing hardware. Independent formula
reconstruction, regression tests, retained hashes, and fresh proof replay
reduce this trust boundary, but neither the encoder nor the complete theorem
has been verified in a proof assistant.

\section{Reproducibility}

The repository contains the encoder, elementary checker, exhaustive branch
audits, formula auditors, compressed proofs, original logs, fresh replay
tools, and machine-readable certificate manifests. The principal commands
are:
\begin{verbatim}
make test
make verify-certificates
make verify-proofs DRAT_TRIM=/path/to/drat-trim
\end{verbatim}

Generated CNFs are not committed. They are rebuilt deterministically and
must match the retained sizes and SHA-256 values before proof replay.

\small
\raggedright
\begin{thebibliography}{10}

\bibitem{GreenwoodGleason1955}
R. E. Greenwood and A. M. Gleason.
\newblock Combinatorial relations and chromatic graphs.
\newblock {\em Canadian Journal of Mathematics}, 7:1--7, 1955.

\bibitem{Exoo1989}
Geoffrey Exoo.
\newblock A lower bound for $R(5,5)$.
\newblock {\em Journal of Graph Theory}, 13(1):97--98, 1989.

\bibitem{McKayRadziszowski1995}
Brendan D. McKay and Stanislaw P. Radziszowski.
\newblock $R(4,5)=25$.
\newblock {\em Journal of Graph Theory}, 19(3):309--322, 1995.

\bibitem{McKayRadziszowski1997}
Brendan D. McKay and Stanislaw P. Radziszowski.
\newblock Subgraph counting identities and Ramsey numbers.
\newblock {\em Journal of Combinatorial Theory, Series B},
69(2):193--209, 1997.

\bibitem{AngeltveitMcKay2018}
Vigleik Angeltveit and Brendan D. McKay.
\newblock $R(5,5)\le48$.
\newblock {\em Journal of Graph Theory}, 89(1):5--13, 2018.

\bibitem{AngeltveitMcKay2025}
Vigleik Angeltveit and Brendan D. McKay.
\newblock $R(5,5)\le46$.
\newblock \href{https://arxiv.org/abs/2409.15709}{arXiv:2409.15709},
version 2, 2025.

\bibitem{Wustep2026}
Wustep.
\newblock Ramsey $R(5,5)$ research record.
\newblock
\href{https://github.com/wustep/maths/tree/4b7b8275d41682a98add180ee05540be84037100/problems/ramsey-r55}
{wustep/maths, problems/ramsey-r55},
audited commit \texttt{4b7b8275d41682a98add180ee05540be84037100},
2026.

\bibitem{Kriebel2026}
Alec Kriebel.
\newblock Prescribed-automorphism cycle-type coverage audit.
\newblock
\href{https://github.com/AlecKriebel/Math/tree/137ffa9f1a340f621651395ad0236cf1bdadb51c/ramsey55}
{AlecKriebel/Math, ramsey55},
audited commit \texttt{137ffa9f1a340f621651395ad0236cf1bdadb51c},
2026.

\end{thebibliography}

\end{document}
